\chapter{Metodologias}

% TODO: Rewrite this thing, Jesus
De forma a prosseguir e conseguir viabilizar construir um software com
qualidade, é necessário adotar algumas metodologias para fundamentar a ideia do
projeto em si, como primeiro passo. Isso ajuda tanto em como priorizar
funcionalidades quanto no aspecto de conseguir gerar dados para construir o
mesmo de forma que entregue o que de fato é necessário para utilização do
mesmo. Com esses elementos em mente, esse capítulo se debruça tanto na
abordagem da metodologia de pesquisa do projeto quanto detalhar o método de
desenvolvimento de software utilizado no projeto.

\section{Metodologia de Desenvolvimento de Software}

Existem diversos métodos de desenvolvimento, cada um com seus respectivos prós
e contras. Todos tentam, de sua forma, responder à necessidade de dar forma ao
processo de desenvolvimento, de forma que ele tente fluir e assegurar que os
requisitos daquele software sejam atendidos, que os membros envolvidos no
processo, tanto de desenvolvimento em si, quanto dos atores envolvidos e
interessados no sistema tenham seus anseios respondidos de forma minimamente
transparente e satisfatória.

Além disso, a escolha de uma boa metodologia, adequada para as necessidades do
sistema, permite que os prazos de entrega sejam melhor controlados e mais
precisos \cite{pressman:2021}, minimizando a ocorrência de falhas e gestão do
tempo.

As ideias fundamentais que norteiam esse projeto são baseadas em dois conceitos
chaves da engenharia de software: desenvolvimento incremental e desenvolvimento
de software ágil. O IEEE define que a engenharia de software é:

\begin{quote}
Engenharia de software: A aplicação de uma abordagem sistemática, disciplinada
e quantificável no desenvolvimento, na operação e na manutenção de software;
isto é, a aplicação de engenharia ao software. (INSTITUTE OF ELECTRICAL AND
ELECTRONICS ENGINEERS, 1990, p. 421, tradução do autor)
\end{quote}

Portanto, a engenharia de software propõe estruturar a criação de software de
forma mais organizada e eficiente. Pressman vai destacar os modelos de
desenvolvimento de software sequenciais e incrementais que moldam boa parte das
práticas de engenharia de software nas últimas décadas.

O modelo em cascata (waterfall) é um modelo sequencial, dos mais tradicionais e
antigos a existirem, cujo qual necessita que cada etapa – levantamento de
requisitos, projeto, implementação, testes – estejam concluídos antes da
próxima etapa começar. Embora forneça uma estrutura rígida e lógica, é
frequentemente criticado pela baixa adaptabilidade em relação às mudanças
durante o processo de desenvolvimento.

Em outro âmbito, o modelo incremental, trabalha com a premissa de um
desenvolvimento feito em pequenas partes, que se agrupam para totalizar o
projeto em si, trabalhando com a concepção de incrementos funcionais, cujo qual
permite maior flexibilidade com as mudanças no processo. Também ajuda a iterar
rapidamente em cima do desenvolvimento do sistema, criando um produto mínimo
viável que possa preencher a proposta usando o mínimo de tempo, reduzindo as
funcionalidades ao seu estado basal. Essa metodologia também ajuda a reduzir
custos e riscos, já que necessita que as funcionalidades estejam em seu estado
mais reduzido, porém funcionais.

Outros modelos tradicionais citados por Pressman incluem o modelo de
prototipação e o modelo em espiral, o primeiro trabalhando na ideia de
elaboração de protótipos pequenos que validem a hipótese, e depois passe para a
implementação final; e o modelo em espiral une a natureza iterativa da
prototipação com os aspectos sistemáticos do modelo em cascata, sendo que acaba
por ser um modelo geralmente usado para desenvolvimento de aplicações em larga
escala, visto que oferece mais segurança ao processo de desenvolvimento. Apesar
da influência histórica, tais modelos acabam por serem um tanto quanto
burocráticos e em um projeto com um escopo de equipe reduzida e onde necessita
um nível de flexibilidade maior, como este, portanto, uma adoção de uma
metodologia mais flexível se faz necessária.

Pressman, comenta que, por meados de 2001, um grupo de desenvolvedores
deflagrou o “Manifesto Ágil”, um documento desenhando uma nova metodologia de
desenvolvimento, destacando a satisfação dos patrocinadores do sistema,
entregas mais rápidas, flexibilidade durante o desenvolvimento do software de
forma a evitar que os desenvolvedores ficaram empacados e reduzisse rigidez. Os
principais destaques da abordagem ágil é a proposta de que a mesma não é
escrita em pedra, ou seja, permite flexibilidade no como aplicá-la. Nesse
cenário, existem algumas abordagens sobre ela, como o Kanban e o Scrum, que dão
um pouco mais de forma a ideia apresentada com o desenvolvimento ágil o que
casa relativamente bem com o modelo de desenvolvimento incremental.

Com base nesse panorama, o presente trabalho adota uma abordagem baseada na
filosofia ágil, estruturada por meio da utilização de um quadro de Kanban,
adaptado à realidade de um desenvolvimento individual. Diferentemente de
modelos tradicionais que envolvem um maior formalismo, o Kanban permite
visualizar e gerenciar o fluxo de trabalho de forma contínua, com foco na
entrega incremental de funcionalidades. Por ser um método não prescritivo, o
Kanban oferece liberdade para que as tarefas possam ser priorizadas, ajustadas
ou redefinidas conforme a evolução do desenvolvimento, característica
especialmente útil quando há apenas um desenvolvedor responsável por todas as
etapas do processo. Vale ressaltar também, que a metodologia adotada acaba por
ajudar a manter o fluxo de desenvolvimento contínuo sem sobrecarga, já que
limita a quantidade de tarefas em progresso, gerando um movimento de maior
produtividade e entregas.


\section{Metodologia de Pesquisa}

\lipsum[1-2]

\subsection{Tipo de Pesquisa}

\lipsum[1]

\subsection{Instrumento de coleta de dados}

\lipsum[1]

\subsection{Análise de Similares}

\lipsum[1]

\subsection{Relação entre os dados coletados e o desenvolvimento}

\lipsum[1]
